% !PRINTED 204
\weberpage{204}
%<<CONT

zur Folge hat, \gesperrt{die Substitution $S_{\nu,\nu'}$ gestattet.} Die Gruppe
der $S_{\nu,\nu'}$ ist aber wegen

\[
S_{\nu,\nu'} S_{\mu,\mu'} = S_{\nu+\mu,\nu'+\mu'}
\]

eine Abelsche, welche nach der Formel

\[
S_{\nu,\nu'} = S_{1,0}^{\nu} S_{0,1}^{\nu'}
\]

durch die Basis $S_{1,0}$, $S_{0,1}$ darstellbar ist, und also ist die \gesperrt{allgemeine Teilungsgleichung eine Abelsche Gleichung und mithin algebraisch lösbar} (Bd. I, § 169\,f.).

\begin{center}
\S\ 61. \quad \gesperrt{Die Teilung der Perioden.}
\end{center}

Die weiter noch zu lösende Aufgabe besteht nun darin, auf
algebraischem Wege die Größen

\weq{1}{x_{\mu,\mu'} = \operatorname{sn}\left(\frac{4\mu K + 4\mu' iK'}{n}\right)}

zu bestimmen. Man erhält alle Werte dieser Größe, wenn man
$\mu$, $\mu'$, voneinander unabhängig, je ein vollständiges Restsystem
nach dem Modul $n$ durchlaufen läßt. Diese Werte sind aber
auch alle \gesperrt{voneinander verschieden}, denn $\operatorname{sn} v$ kann nur dann
$= \operatorname{sn} v'$ sein, wenn $v'$ kongruent $v$ oder kongruent $2K - v$
modulo $4K$, $2iK'$, und beides kann für zwei verschiedene der
Argumente von (1) nicht eintreten.

Die $n^2$ Größen (1) sind die Wurzeln der Gleichung

\weq{2}{xA(x^2) = 0,}

und nach den Formeln

\weq{3}{y = \frac{D(x^2)}{B(x^2)}, \qquad z = \frac{D(x^2)}{C(x^2)}}

können auch

\[
\begin{aligned}
y_{\mu,\mu'} &= \operatorname{cn}\left(\frac{4\mu K + 4\mu' iK'}{n}\right),\\
z_{\mu,\mu'} &= \operatorname{dn}\left(\frac{4\mu K + 4\mu' iK'}{n}\right)
\end{aligned}
\]

rational durch $x_{\mu,\mu'}$ ausgedrückt werden, wenn der \gesperrt{Rationalitätsbereich aus rationalen Zahlen und rationalen
Funktionen von $x^2$ besteht}.

Aus den Wurzeln $x^2$ der Gleichung $A = 0$ lassen sich nach
§ 60, (2) und § 44, (18), (19), (20) die Wurzeln von $B = 0$,
$C = 0$, $D = 0$ rational ableiten; wenn nämlich $x^2$ eine Wurzel
von $A = 0$ ist, so sind
%CONT>>
